\chapter{Gần và Xa}
\markboth{Gần và Xa}{Near and Far}

\section{Ở gần nhau nhưng ai cũng cắm mặt vào điện thoại.}
\begin{truyen}{Ở gần nhau nhưng ai cũng cắm mặt vào điện thoại.}{Near bodies can hide far hearts.}
\chuhoa{B}{ữa cơm có đủ người, nhưng câu chuyện không nối thành câu. Mỗi người chạm vào một màn hình riêng, thỉnh thoảng ngẩng lên vài giây rồi lại trôi đi. Khoảng cách lớn nhất đôi khi không nằm ở số cây số, mà nằm ở chỗ không còn ai thật sự hiện diện.}
\textit{(Everyone was physically present at the meal, yet no real conversation formed. Each person touched a separate screen, looking up only for a few seconds before drifting away again. The greatest distance is sometimes not measured in kilometers, but in the absence of real presence.)}
\end{truyen}
\begin{baihoc}
Gần về chỗ ngồi chưa chắc gần về tâm hồn.
\textit{(Sitting near does not guarantee being inwardly close.)}
\end{baihoc}
\begin{ghinhoanh}
\textbf{Short English to remember:}

Presence is more than distance.
\end{ghinhoanh}
\begin{tuhocnhanh}
1. Khi ở cạnh người thân, em có thật sự ở đó không? \textit{(When you are beside loved ones, are you really there?)}

2. Điều gì đang làm những cuộc gặp của em thành hời hợt? \textit{(What is making your meetings become shallow?)}
\end{tuhocnhanh}
\ngancach

\section{Xa nhau một thời gian để hiểu nhau rõ hơn.}
\begin{truyen}{Xa nhau một thời gian để hiểu nhau rõ hơn.}{Distance can reveal value.}
\chuhoa{C}{ó những ngày em ở nội trú, không còn nghe tiếng mẹ nhắc dậy học mỗi sáng. Chính quãng xa ấy làm em nhận ra những điều từng thấy phiền thực ra là chăm sóc. Khi một người tạm ra khỏi tầm tay, lòng biết ơn nhiều khi mới bước vào tầm nhìn.}
\textit{(There were boarding-school days when the student no longer heard her mother waking her for school. That period of distance revealed that what once seemed annoying was actually care. When someone moves out of our reach, gratitude often moves into our sight.)}
\end{truyen}
\begin{baihoc}
Xa không chỉ là mất kết nối; đôi khi là cách giá trị hiện hình.
\textit{(Distance is not only disconnection; sometimes it is how value becomes visible.)}
\end{baihoc}
\begin{ghinhoanh}
\textbf{Short English to remember:}

Distance can uncover love.
\end{ghinhoanh}
\begin{tuhocnhanh}
1. Có ai em chỉ thật sự quý khi họ không còn ở gần mỗi ngày? \textit{(Is there anyone you only truly appreciated when they were no longer near every day?)}

2. Em có đang quen quá nên quên trân trọng không? \textit{(Have you become so used to someone that you forgot to appreciate them?)}
\end{tuhocnhanh}
\ngancach

\section{Một người quá gần với lời khen nên không nghe được góp ý.}
\begin{truyen}{Một người quá gần với lời khen nên không nghe được góp ý.}{Too close to praise, too far from truth.}
\chuhoa{K}{hi quanh mình toàn lời tán thưởng, anh bắt đầu khó chịu với mọi góp ý trái tai. Anh không nhận ra mình đang đứng rất gần cảm giác dễ chịu, nhưng lại ngày càng xa sự thật cần thiết cho trưởng thành.}
\textit{(Surrounded by praise, he gradually became irritated by any uncomfortable feedback. He did not realize that he had moved very close to comfort while drifting far away from the truth needed for growth.)}
\end{truyen}
\begin{baihoc}
Điều ở gần tai chưa chắc là điều ở gần chân lý.
\textit{(What stays near your ear is not always near the truth.)}
\end{baihoc}
\begin{ghinhoanh}
\textbf{Short English to remember:}

Comfort can mute truth.
\end{ghinhoanh}
\begin{tuhocnhanh}
1. Em có đang chỉ thích nghe điều dễ chịu không? \textit{(Are you only willing to hear what feels pleasant?)}

2. Góp ý nào em nên nghe lại bằng thái độ bình tĩnh hơn? \textit{(Which feedback should you revisit more calmly?)}
\end{tuhocnhanh}
\ngancach

\section{Một học sinh ngồi cuối lớp nhưng lại rất gần việc học.}
\begin{truyen}{Một học sinh ngồi cuối lớp nhưng lại rất gần việc học.}{Where you sit is not where you stand within.}
\chuhoa{E}{m không ngồi bàn đầu, cũng ít nói trước lớp. Nhưng em ghi chép kỹ, hỏi riêng khi chưa hiểu, và tự học đều mỗi tối. Có những người ở vị trí không nổi bật mà lại đứng rất gần con đường đi lên của chính mình.}
\textit{(The student did not sit in the front row and rarely spoke publicly. Yet the notes were careful, questions were asked privately, and self-study continued every night. Some people occupy unremarkable positions while standing very near their own path upward.)}
\end{truyen}
\begin{baihoc}
Sự gần gũi với mục tiêu không phải lúc nào cũng lộ ra bên ngoài.
\textit{(Closeness to a goal is not always outwardly visible.)}
\end{baihoc}
\begin{ghinhoanh}
\textbf{Short English to remember:}

Quiet effort can stay close to growth.
\end{ghinhoanh}
\begin{tuhocnhanh}
1. Em có đang đánh giá người khác quá nhanh theo vẻ ngoài không? \textit{(Are you judging others too quickly by appearance?)}

2. Em đang gần ước mơ của mình ở điểm nào mà người khác không thấy? \textit{(In what way are you close to your dream even if others cannot see it?)}
\end{tuhocnhanh}
\ngancach

\section{Giữ một khoảng xa vừa đủ để tôn trọng nhau.}
\begin{truyen}{Giữ một khoảng xa vừa đủ để tôn trọng nhau.}{Healthy distance protects respect.}
\chuhoa{K}{hông phải thân là được phép chen vào mọi quyết định của nhau. Một người chị thương em mình nhưng không kiểm soát hộ em từng việc. Nhờ vậy, tình cảm không biến thành ngột ngạt, và khoảng xa vừa đủ lại giúp mối quan hệ đứng lâu hơn.}
\textit{(Closeness does not give us permission to enter every decision of another person. An older sister loved her younger sibling but did not control every step. Because of that, affection did not become suffocation, and a healthy distance helped the relationship endure.)}
\end{truyen}
\begin{baihoc}
Xa vừa đủ không làm tình cảm yếu đi; nó giúp tình cảm thở được.
\textit{(Proper distance does not weaken affection; it helps affection breathe.)}
\end{baihoc}
\begin{ghinhoanh}
\textbf{Short English to remember:}

Respect needs space.
\end{ghinhoanh}
\begin{tuhocnhanh}
1. Em có đang thương ai đó theo cách quá kiểm soát không? \textit{(Are you caring for someone in an overly controlling way?)}

2. Mối quan hệ nào của em đang cần thêm không gian? \textit{(Which relationship in your life needs a little more space?)}
\end{tuhocnhanh}
\ngancach
